\documentclass[11pt,a4paper]{article}
\usepackage[italian]{babel}
\usepackage[margin=2.5cm]{geometry}
\usepackage{parskip}
\usepackage{amsmath, amssymb, amsthm}
\usepackage{booktabs}
\usepackage{array}
\usepackage{xcolor}
\usepackage{needspace}
\usepackage{enumitem}
\usepackage{listings}
\usepackage{tikz}
\usepackage[most]{tcolorbox}
\usepackage[hidelinks]{hyperref}
\usetikzlibrary{decorations.pathreplacing, shapes.geometric, arrows.meta, positioning}

% Niente vedove/orfane: i paragrafi non lasciano righe isolate a fine/inizio pagina
\widowpenalty=10000
\clubpenalty=10000

% Elenchi più compatti
\setlist{itemsep=2pt, parsep=0pt, topsep=4pt, partopsep=0pt}

% --- Riquadri con bordo nero, senza numero (si spezzano tra due pagine) ---
% Uso: \begin{definizione} ... \end{definizione}
%      \begin{teorema}[Nome] ... \end{teorema}  (il nome è facoltativo)
\tcbset{nota/.style={enhanced jigsaw, breakable, colback=white,
  colframe=black, boxrule=0.5pt, sharp corners}}
\NewTColorBox{definizione}{}{nota,
  before upper={\textbf{Definizione.}\ }}
\NewTColorBox{teorema}{o}{nota,
  before upper={\textbf{Teorema\IfValueT{#1}{ (#1)}.}\ }}
\NewTColorBox{proposizione}{o}{nota,
  before upper={\textbf{Proposizione\IfValueT{#1}{ (#1)}.}\ }}
\NewTColorBox{proprieta}{o}{nota,
  before upper={\textbf{Proprietà\IfValueT{#1}{ (#1)}.}\ }}
\NewTColorBox{assioma}{o}{nota, boxrule=1pt,
  before upper={\textbf{Assioma\IfValueT{#1}{ (#1)}.}\ }}

% --- Ambienti semplici (senza riquadro) ---
% Il titolo sta in cima al blocco, anche se il contenuto inizia con un riquadro o
% con codice; e non resta mai da solo in fondo a una pagina.
% Uso: \begin{esempio}[Titolo facoltativo] ... \end{esempio}
\newcommand{\newrunin}[2]{%
  \NewDocumentEnvironment{#1}{o}{%
    \par\Needspace{4\baselineskip}\addvspace{\medskipamount}\noindent
    \textbf{#2}\IfValueT{##1}{ (##1)}\textbf{.}\ \ignorespaces}%
  {\par\addvspace{\medskipamount}}}
\newrunin{esempio}{Esempio}
\newrunin{esercizio}{Esercizio}
\newrunin{osservazione}{Osservazione}

% V e F nelle tavole di verità
\newcommand{\V}{\textsf{V}}
\newcommand{\F}{\textsf{F}}

% --- Codice C ---
% Uso: \begin{codice} ... \end{codice}      codice C numerato
%      \begin{comando} ... \end{comando}    comandi da terminale
%      \begin{risultato} ... \end{risultato} output di un programma
%      \lstinline|printf("ciao");|           codice dentro al testo
\lstdefinestyle{c}{
  language=C,
  basicstyle=\ttfamily\small,
  keywordstyle=\color{blue}\bfseries,
  commentstyle=\color{gray}\itshape,
  stringstyle=\color{red!70!black},
  numbers=left,
  numberstyle=\tiny\color{gray},
  numbersep=8pt,
  columns=flexible,
  keepspaces=true,
  breaklines=true,
  showstringspaces=false,
  tabsize=4,
  morekeywords={include, define},
  % lettere accentate nei commenti
  literate={à}{{\`a}}1 {è}{{\`e}}1 {é}{{\'e}}1 {ì}{{\`i}}1
           {ò}{{\`o}}1 {ù}{{\`u}}1 {È}{{\`E}}1 {À}{{\`A}}1
}
\lstset{style=c}
\newtcblisting{codice}{listing only, nota, left=6mm,
  listing options={style=c}}
\newtcblisting{comando}{listing only, nota,
  listing options={basicstyle=\ttfamily\small, numbers=none, language={},
    columns=flexible, keepspaces=true, breaklines=true}}
\newtcblisting{risultato}{listing only, enhanced jigsaw, breakable,
  colback=black!5, colframe=black, boxrule=0.5pt, sharp corners,
  listing options={basicstyle=\ttfamily\small, numbers=none, language={},
    columns=flexible, keepspaces=true, breaklines=true}}

% --- Diagrammi di flusso (simboli standard) ---
\tikzset{
  terminale/.style = {ellipse, draw, minimum width=2.5cm, minimum height=0.9cm},
  io/.style        = {trapezium, trapezium left angle=70, trapezium right angle=110,
                      draw, minimum width=2.5cm, minimum height=0.9cm},
  processo/.style  = {rectangle, draw, minimum width=2.5cm, minimum height=0.9cm},
  decisione/.style = {diamond, draw, aspect=2, minimum width=2.5cm},
  freccia/.style   = {thick, -{Stealth}}
}

\title{Funzioni esponenziali, logaritmiche e trigonometriche}
\author{Marco Cerbella}
\date{Lezione 7 -- 15 ottobre 2026}

\begin{document}
\maketitle

\section{Funzioni esponenziali}

\begin{definizione}[Funzione esponenziale]
Dato $a \in \mathbb{R}$, $a > 0$, $a \ne 1$, si chiama \emph{funzione esponenziale in base $a$} la funzione
\[
f(x) = a^x .
\]
Il dominio è $\mathbb{R}$ e l'immagine è $(0, +\infty)$.
\end{definizione}

Se $a = 1$ si avrebbe $f(x) = 1$, che è la funzione costante. Valgono le regole delle potenze:
\[
a^0 = 1, \qquad a^x > 0 \ \ \forall x \in \mathbb{R}, \qquad a^{x+y} = a^x \cdot a^y, \qquad (a^x)^y = a^{xy}.
\]

La funzione è \emph{strettamente crescente} se $a > 1$ e \emph{strettamente decrescente} se $a < 1$. È limitata inferiormente (da $0$) ma non superiormente.
\begin{itemize}
    \item Se $a > 1$: al crescere di $x$, $a^x$ diventa molto grande; per $x$ negativo grande in valore assoluto, $a^x = \dfrac{1}{a^{-x}}$ si avvicina a $0$.
    \item Se $a < 1$: si ha $a^x = \left(\dfrac{1}{b}\right)^x = b^{-x}$ con $b = \dfrac{1}{a} > 1$. Il grafico di $a^x$ si ottiene dal grafico di $b^x$ con una \emph{riflessione rispetto all'asse $y$}.
\end{itemize}

\begin{esempio}[Test]
$f(x) = e^{x+1}$ è una funzione: (1) limitata superiormente ma non inferiormente, (2) limitata inferiormente ma non superiormente, (3) pari, (4) dispari. La risposta corretta è la (2): $e^{x+1} > 0$ ma cresce senza limite.
\end{esempio}

\section{Funzioni logaritmiche}

\begin{definizione}[Logaritmo]
Per $a > 0$, $a \ne 1$, $b > 0$ e $c \in \mathbb{R}$,
\[
\log_a b = c \iff a^c = b .
\]
Il numero $a$ è la \emph{base} del logaritmo.
\end{definizione}

\begin{definizione}[Funzione logaritmica]
Dato $a > 0$, $a \ne 1$, la \emph{funzione logaritmica in base $a$} è $f(x) = \log_a x$. Il dominio è $(0, +\infty)$ e l'immagine è $\mathbb{R}$.
\end{definizione}

La funzione è strettamente crescente se $a > 1$ e strettamente decrescente se $0 < a < 1$. Per $a > 1$, $\log_a x$ diventa molto grande per $x$ molto grande e molto negativo (tende a $-\infty$) per $x$ che si avvicina a $0$; per $0 < a < 1$ succede il contrario.

\begin{proprieta}
Per $a > 0$, $a \ne 1$:
\begin{itemize}
    \item $\log_a 1 = 0$ e $\log_a a = 1$;
    \item $\log_a (xy) = \log_a x + \log_a y$ per $x, y > 0$;
    \item $\log_a x^\alpha = \alpha \log_a x$ per $\alpha \in \mathbb{R}$, $x > 0$;
    \item (cambio di base) $\log_b x = \dfrac{\log_a x}{\log_a b}$ per $b > 0$, $b \ne 1$, $x > 0$.
\end{itemize}
\end{proprieta}

\subsection*{Relazione tra esponenziale e logaritmo}

La funzione esponenziale $f : \mathbb{R} \to (0,+\infty)$, $f(x) = a^x$, è monotona, quindi invertibile. La funzione logaritmo $g : (0,+\infty) \to \mathbb{R}$, $g(x) = \log_a x$, è la sua \emph{inversa}: $g = f^{-1}$.

\section{Funzioni trigonometriche}

\subsection{La circonferenza goniometrica}

La \emph{circonferenza goniometrica} è la circonferenza di centro $O$ e raggio $1$. Serve a rappresentare tutti gli angoli, da quello nullo ($0^\circ$) all'angolo giro ($360^\circ$): ogni angolo al centro individua un arco, e viceversa. Il punto di partenza è $A = (1, 0)$.

Gli angoli si misurano in gradi o in \emph{radianti}: un radiante è l'angolo al centro che sottende un arco di lunghezza uguale al raggio, quindi l'angolo giro misura $2\pi$ radianti. Gli angoli positivi si ottengono percorrendo la circonferenza in senso antiorario a partire da $A$, quelli negativi in senso orario. Gli angoli oltre $2\pi$ ripassano sugli stessi punti: l'angolo $0$ coincide con $2\pi$.

\begin{center}
\renewcommand{\arraystretch}{1.2}
\begin{tabular}{@{}l|cccccccc@{}}
\toprule
$\theta$ (gradi) & $0$ & $30$ & $45$ & $60$ & $90$ & $180$ & $270$ & $360$ \\
$\theta$ (radianti) & $0$ & $\frac{\pi}{6}$ & $\frac{\pi}{4}$ & $\frac{\pi}{3}$ & $\frac{\pi}{2}$ & $\pi$ & $\frac{3\pi}{2}$ & $2\pi$ \\
\bottomrule
\end{tabular}
\end{center}

\subsection{Seno, coseno, tangente, cotangente}

Sia $P = (x_P, y_P)$ il punto della circonferenza che corrisponde all'angolo $\theta$.
\begin{itemize}
    \item $\sin\theta$ è l'\emph{ordinata} di $P$ (il segmento $PH$ con segno);
    \item $\cos\theta$ è l'\emph{ascissa} di $P$ (il segmento $OH$ con segno);
    \item $\tan\theta$ è l'ordinata del punto di intersezione tra la retta parallela all'asse $y$ passante per $A$ e la retta che unisce $O$ a $P$. Non è definita per $P = (0, \pm 1)$, cioè $\theta = \frac{\pi}{2} + k\pi$;
    \item $\cot\theta$ è l'ascissa del punto di intersezione tra la retta parallela all'asse $x$ passante per $M = (0, 1)$ e la retta $OP$. Non è definita per $P = (\pm 1, 0)$, cioè $\theta = k\pi$.
\end{itemize}

\subsection{Le funzioni}

\begin{center}
\renewcommand{\arraystretch}{1.3}
\begin{tabular}{@{}llllc@{}}
\toprule
Funzione & Dominio & Immagine & Periodo & Limitata \\
\midrule
$\sin x$ & $\mathbb{R}$ & $[-1, 1]$ & $2\pi$ & sì \\
$\cos x$ & $\mathbb{R}$ & $[-1, 1]$ & $2\pi$ & sì \\
$\tan x$ & $\mathbb{R} \setminus \{\frac{\pi}{2} + k\pi,\ k \in \mathbb{Z}\}$ & $\mathbb{R}$ & $\pi$ & no \\
$\cot x$ & $\mathbb{R} \setminus \{k\pi,\ k \in \mathbb{Z}\}$ & $\mathbb{R}$ & $\pi$ & no \\
\bottomrule
\end{tabular}
\end{center}

Il seno è dispari e il coseno è pari.

\subsection{Relazioni trigonometriche fondamentali}

Per ogni $x, x_1, x_2 \in \mathbb{R}$:
\begin{itemize}
    \item $\cos^2 x + \sin^2 x = 1$ (teorema di Pitagora nel triangolo rettangolo $OPH$, con $OP = 1$);
    \item $\cos(x_1 + x_2) = \cos x_1 \cos x_2 - \sin x_1 \sin x_2$;
    \item $\sin(x_1 + x_2) = \sin x_1 \cos x_2 + \cos x_1 \sin x_2$.
\end{itemize}
Dalle formule di somma derivano
\[
\cos(-x) = \cos x, \quad \sin(-x) = -\sin x, \quad \cos 2x = \cos^2 x - \sin^2 x, \quad \sin 2x = 2\sin x \cos x.
\]
Dalle definizioni (triangoli simili) si ha inoltre
\[
\tan x = \frac{\sin x}{\cos x}, \qquad \cot x = \frac{\cos x}{\sin x}.
\]

\subsection{Funzioni trigonometriche inverse}

Le funzioni trigonometriche non sono invertibili su tutto il dominio perché sono periodiche. Si restringe il dominio a un intervallo su cui sono strettamente monotone.
\begin{center}
\renewcommand{\arraystretch}{1.4}
\begin{tabular}{@{}lll@{}}
\toprule
Restrizione & Funzione inversa & Dominio dell'inversa \\
\midrule
$y = \sin x$, $x \in \left[-\frac{\pi}{2}, \frac{\pi}{2}\right]$ & $x = \arcsin y$ & $y \in [-1, 1]$ \\
$y = \cos x$, $x \in [0, \pi]$ & $x = \arccos y$ & $y \in [-1, 1]$ \\
$y = \tan x$, $x \in \left(-\frac{\pi}{2}, \frac{\pi}{2}\right)$ & $x = \arctan y$ & $y \in \mathbb{R}$ \\
$y = \cot x$, $x \in (0, \pi)$ & $x = \operatorname{arccot} y$ & $y \in \mathbb{R}$ \\
\bottomrule
\end{tabular}
\end{center}

\begin{esempio}[Test]
$f(x) = \sin x + \cos x$ è una funzione: (1) limitata, (2) invertibile, (3) pari, (4) dispari. La risposta corretta è la (1): $|\sin x + \cos x| \le |\sin x| + |\cos x| \le 2$. Non è invertibile perché periodica e non è né pari né dispari, dato che $f(-x) = -\sin x + \cos x$.
\end{esempio}

\end{document}
